\scanpage{28}
\margin{\uline{à vérifier}}

Nous laissons ici les détails de la vérification de l'existence de $Z',\nu'$.

On prouve (ii) pr \deleted{\unsure{la|les}} \inserted{$=$} méthode que (i), \deleted{\unsure{à nouveau}} en se ramenant au cas où $z,z'$ sont définis en termes des $=$ $T,\nu$ et deux cycles $Z,Z'$ sur $T$. On \inserted{peut} se ramener aussi à (iii) et à l'énoncé analogue à (ii) pr les produits \uline{extérieurs}, ce qui dispense de faire encore bouger les cycles.

La démonstration de (iv) se fait pr l'\unsure{AQT}. Celle de (v) est triviale sur les définitions (ou sur le formalisme de l'équiv. alg.)

\margin{\figure{\begin{tikzpicture}\draw (0,0)--(0,3);\end{tikzpicture}}}

\deleted{\uline{Théorème}} \uline{\unsure{Scholie|Scolie}} 1.3. Les $A'^*(X)$, pr $X$ variable, ont les $=$s propriétés formelles générales que les $A^*(X)$. En particulier, on a la formule de projection, Dans les cas où équiv. linéaire $=$ équiv. Alg. sur $X$ (p. ex. \unsure{l'espace} ponctuel, les $\mathbb P^n$, les variétés de drapeaux) \unsure{ils sont} $=$ équiv. alb.

\mspage{27}
